% =========================================================================
\section{Ownership, Routing and Rules (RQ3)}
\label{sec:layers}

\subsection{Ownership and routing are weakly related}
\label{sec:wealth-power}
Both ledgers are extremely unequal in ownership---balance Gini
\V{link.gini.all} and \V{uni.gini.all}, with the top 1\% of addresses holding
\V{link.top1} and \V{uni.top1} of supply---yet unconcentrated in routing,
with a mean weekly flow HHI of \V{link.fhhi.weekly} and \V{uni.fhhi.weekly}.
Across the addresses of the $10^{4}$ subgraph that hold a positive balance,
the rank correlation between balance and counterparty count is
\V{link.rho.cp.pos} on LINK ($n = \V{link.rho.n.pos}$, 95\% CI
\V{link.rho.cp.pos.ci}) and \V{uni.rho.cp.pos} on UNI
($n = \V{uni.rho.n.pos}$, \V{uni.rho.cp.pos.ci}); between balance and
betweenness it is \V{link.rho.bet.pos} (\V{link.rho.bet.pos.ci}) and
\V{uni.rho.bet.pos} (\V{uni.rho.bet.pos.ci}). Including the addresses with
zero balance gives \V{link.rho.cp.all}
and \V{uni.rho.cp.all} for counterparties. These correlations are small in
magnitude; with samples of this size some are nonetheless distinguishable
from zero, and we read them as weak association rather than independence.
Weak association is what allows extreme ownership inequality to coexist with
an unconcentrated routing layer: knowing who holds the supply says little
about who moves it.

\subsection{The rule layer, measured in delegated votes}
\label{sec:governance}
Uniswap governance is a threshold game. Submitting an on-chain proposal
requires $10^{6}$ UNI of votes delegated to the proposer, and a proposal
passes with $4\times10^{7}$ UNI of votes in favour
\cite{uniswap2026governance}; we read both values from the governor contract
at the first and last block of the window. Votes are not balances. A token
carries voting weight only once its holder delegates it, to itself or to
another address, and the weight is checkpointed per block. We therefore
measure the rule layer in delegated votes at the snapshot block, read from
the token contract: we query the delegate of every one of the \V{uni.holders}
holders and the checkpointed votes of every one of the
\V{uni.gov.n_delegates} delegates, and verify that each delegate's votes
equal the summed balance of the holders delegating to it.

\paragraph{Specification again}
Table~\ref{tab:gov} sets the delegated-vote measurement beside the
balance-based one used in the first version of this paper.
\V{uni.gov.deleg_share} of the supply was delegated at the snapshot. Counting
balances, the smallest set of voting-capable addresses holding $4\times10^{7}$
UNI has \V{uni.gov.bal.nonvoting.n_quorum} members; the two addresses in
question are a Binance hot wallet and an unlabelled address
(\V{uni.gov.nonvoting.m2}), and neither has delegated, so their voting weight
is zero. Excluding custodial addresses as well removes the Binance wallet;
its place is taken by a second unlabelled address
(\V{uni.gov.custodial.m2}), which has not delegated either. Counting delegated votes,
\V{uni.gov.n_quorum} delegates are needed to reach the quorum and
\V{uni.gov.n_half} to hold a majority of all delegated votes---close to the
eleven that Fritsch et al.\ report for Uniswap in March 2022
\cite{fritsch2024voting}. \V{uni.gov.n_eligible} delegates individually hold
enough votes to submit a proposal. A balance-based count thus overstates
concentration in this instance, and the direction of the error is not
knowable in advance: it depends on who delegates.

\begin{table}[t]
\centering\small
\caption{The UNI rule layer under four specifications at block
\V{snapshot.block}. $N_{40\mathrm{M}}$: fewest addresses whose combined weight
reaches the $4\times10^{7}$ UNI quorum. $N_{50\%}$: fewest addresses holding
half of the total weight. $\geq 1\mathrm{M}$: addresses that individually
meet the proposal threshold. Only delegated votes carry voting weight.}
\label{tab:gov}
\setlength{\tabcolsep}{3.5pt}
\begin{tabular}{@{}lcccc@{}}
\toprule
Specification & $N_{40\mathrm{M}}$ & $N_{50\%}$ & $\geq 1\mathrm{M}$ & Weight / supply \\
\midrule
Balances, all holders & \V{uni.gov.bal.all.n_quorum} & \V{uni.gov.bal.all.n_half} & \V{uni.gov.bal.all.n_eligible} & \V{uni.gov.bal.all.share} \\
Balances, excl.\ non-voting & \V{uni.gov.bal.nonvoting.n_quorum} & \V{uni.gov.bal.nonvoting.n_half} & \V{uni.gov.bal.nonvoting.n_eligible} & \V{uni.gov.bal.nonvoting.share} \\
\quad and excl.\ custodial & \V{uni.gov.bal.custodial.n_quorum} & \V{uni.gov.bal.custodial.n_half} & \V{uni.gov.bal.custodial.n_eligible} & \V{uni.gov.bal.custodial.share} \\
\textbf{Delegated votes} & \textbf{\V{uni.gov.n_quorum}} & \textbf{\V{uni.gov.n_half}} & \textbf{\V{uni.gov.n_eligible}} & \V{uni.gov.deleg_share} \\
\bottomrule
\end{tabular}
\end{table}

\paragraph{The rule layer has already rewritten the routing layer}
On 25 December 2025 Uniswap governance passed the ``UNIfication'' proposal,
which activated protocol fees on Uniswap v2 and selected v3 pools, routed
those fees to a \texttt{TokenJar} contract on each chain, and burned $10^{8}$
UNI from the treasury; the burn was executed on 28 December
\cite{theblock2025unification, cointelegraph2025burn, uniswap2026fees}.
Assets leave a \texttt{TokenJar} only through the releaser contract it
designates, and the jar's owner---on Ethereum, the governance timelock---can
change that releaser. Throughout our window the Ethereum jar's releaser was the
Firepit contract listed in Uniswap's deployment documentation
\cite{uniswap2026deployments}, which releases the accumulated fees to whoever
burns \V{uni.firepit.threshold} UNI; we read both values from the chain at
the first and last block of the window. Its effects are visible in our data.
All \V{uni.burn.n} UNI records that a naive pipeline would discard as
``burns'' are transfers to \texttt{0x\ldots dEaD}: \V{uni.burn.hi.n} in lots
of \V{uni.burn.lot_hi} UNI from \V{uni.burn.hi.senders} senders on
Ethereum; \V{uni.burn.lo.n} in lots of \V{uni.burn.lot_lo} UNI, of which
\V{uni.burn.lo.bridged} arrive through the Base, Arbitrum, Optimism and
Unichain bridges and \V{uni.burn.lo.other} from an unlabelled address---a lot
size that matches the \V{uni.burn.lot_lo}-UNI threshold we read at the end
of the window from the releasers on \V{uni.l2.chains}, each identified as
the address that the deployment page lists and that the chain's
\texttt{TokenJar} designates on chain (addresses and calls in the
repository, \ref{app:repro}); and one of \V{uni.burn.odd} UNI. They total \V{uni.burn.window} UNI. The burn address holds
\V{uni.burn_pct} of supply at the snapshot, of which $10^{8}$ UNI is the
one-off December burn. A vote of the rule layer changed the economics of
every swap routed through a fee-enabled pool, which is the strongest
evidence we have that routing contestability is held conditionally.

\paragraph{The rule layer and the routing layer do not overlap}
None of the \V{uni.gov.n_eligible} delegates that can submit a proposal is
among the 10, 50 or 100 largest senders of UNI by mean weekly outflow, and
\V{uni.gov.inactive} of them (\V{uni.gov.inactive_pct}) sent or received no
UNI at all during the window. The largest senders have delegated none of
their own balances. Because a hidden broker is defined to hold at most the
median balance, a comparison of brokers with proposal-eligible addresses
would be empty by construction; the outflow-based comparison has no balance
condition and could have found an overlap.

\paragraph{Routing under direct governance control}
Of mean weekly UNI outflow, \V{uni.gov.rev_l1} originates from contracts the
governance timelock controls and \V{uni.gov.rev_l2} from other
Uniswap-labelled contracts; \V{uni.gov.rev_l3} comes from addresses the label
compilation does not identify. This bounds direct control of UNI
\emph{token} routing, not of swap volume in Uniswap pools, which the fee
switch shows governance can reach.

LINK has no token-weighted governance: holding LINK confers no vote, and
while anyone may run a Chainlink oracle, each oracle network can restrict
which oracles contribute to it \cite{chainlink2022nodeoperator}. Of the two
ledgers whose routing layers we measure, only one has a rule layer that
token holders can use to change it.
